\section{Estrategia}

Quien está frente a varias opciones parecidas no necesita otra tabla de nutrientes. Necesita poder interpretar lo que sí hay y compararlo según lo que importa en esa compra. NutriMatch no decide por la persona ni dice cuál alimento es mejor en general. Ayuda a encontrar, entender y comparar opciones para tomar una decisión de compra con la información disponible y con las prioridades que ella declara.

\subsection{Usuarios objetivo}

La aplicación está pensada para quien compra alimentos empacados y, ante productos similares, quiere más elementos para elegir. Puede ser la compra de la semana, la búsqueda de un sustituto de algo que ya consume o simplemente una situación en la que dos opciones parecen iguales.

No hace falta saber de nutrición ni de Ciencia de Datos. Hace falta poder decir, de forma sencilla, qué importa más al elegir: la información nutricional, el nivel de procesamiento o ciertas características del producto. También se puede indicar si la alimentación es vegetariana o vegana y si hay que evitar algún alérgeno.

Esa persona no busca un diagnóstico ni una recomendación médica. Busca decidir la compra con lo que los datos permiten afirmar.

No se observó a compradores reales ni se midieron características como edad, ingreso o ciudad. El perfil que aparece en la evaluación es un ejemplo para observar qué ocurre cuando cambian las prioridades. No representa a una persona registrada.

\subsection{Necesidad}

La necesidad aparece cuando hay varias opciones parecidas y la información para compararlas está repartida en la etiqueta, o no está igual de completa en todos los productos. Revisar a mano nutrientes, sellos, ingredientes, restricciones y precio vuelve más difícil la comparación. Si además falta un dato, leer ese silencio como si la característica no existiera puede cambiar el sentido de la elección.

Lo que hace falta, entonces, es reunir la información disponible y presentarla de modo que la persona pueda pasar de tener varias opciones frente a ella a elegir entre ellas.

\subsection{Accionables}

El resultado no se queda en un número. Sirve para decidir qué hacer con el producto que se está mirando. La persona puede conservarlo, descartarlo, buscar una alternativa, sustituirlo o revisar con más cuidado alguna restricción.

En el carrito, el conjunto elegido también se puede observar mediante los grupos del Plato del Bien Comer: verduras y frutas, cereales, leguminosas y alimentos de origen animal, además de lo que no se puede clasificar \citep{nom0432013}. Esta vista da contexto a la compra y permite revisar el conjunto elegido. No es un diagnóstico.

\begin{table}[ht]
  \centering
  \caption{Acciones de compra que permite el resultado.}
  \label{tab:accionables}
  \small
  \begin{tabularx}{\textwidth}{>{\raggedright\arraybackslash}p{3.4cm}XX}
    \toprule
    Acción & Qué puede hacer la persona & Para qué le sirve \\
    \midrule
    Conservar un producto &
    Revisar el resultado y decidir si lo lleva. &
    Elegir con base en las prioridades que declaró. \\
    \addlinespace
    Descartarlo &
    Dejar fuera una opción que no corresponde a lo que busca o que no se puede verificar. &
    No quedarse con un producto solo porque se parece a los demás. \\
    \addlinespace
    Explorar alternativas &
    Consultar otros productos del mismo tipo. &
    Ver opciones comparables sin recorrer todo el catálogo. \\
    \addlinespace
    Sustituirlo &
    Cambiar el producto por otra opción de ese mismo tipo. &
    Ajustar la elección sin empezar la búsqueda de cero. \\
    \addlinespace
    Revisar restricciones &
    Consultar información sobre dieta y alérgenos. &
    Saber si el producto no cumple una condición o si no hay información suficiente para confirmarlo. \\
    \addlinespace
    Cerrar la compra &
    Llevar lo elegido al carrito y revisar el conjunto. &
    Pasar de la comparación de un producto a la decisión de compra. \\
    \bottomrule
  \end{tabularx}
\end{table}

Estos accionables conectan el resultado del sistema con una decisión concreta. NutriMatch no busca que la persona interprete un puntaje por sí solo, sino que pueda utilizar la comparación para decidir qué producto conservar, cuál descartar o por cuál sustituirlo.

\subsection{Propuesta de solución}

La comparación empieza con una pregunta sencilla: ¿qué es más importante en esta compra?

La persona ordena nutrición, procesamiento y las características que valora. A partir de ese orden se organizan los productos y se muestra cuáles se ajustan más a esas prioridades.

Se puede empezar por un producto concreto, ver sus características y después mirar alternativas del mismo tipo. No hace falta recorrer el catálogo entero para encontrar un sustituto.

El resultado se lee sin interpretar una fórmula. Se distingue si el producto entra en la comparación, si la información no alcanza o si una condición no se puede confirmar. El precio se muestra cuando existe una observación real. Si el monto es de demostración, se identifica como tal: informa la ficha y el recorrido de compra, pero no decide el orden.

En resumen, no se le dice a la persona qué debe comer. Se convierte información dispersa en una comparación que puede utilizar en el momento de comprar. La forma en que esa comparación se construye a partir de los datos se explica en las secciones siguientes.

\subsection{Usabilidad}

El recorrido sigue el orden de una compra: identificar un producto, entenderlo, comparar alternativas y decidir si se lleva.

No hace falta conocer el cálculo. Se busca un producto, se revisa su ficha, se consultan alternativas y se elige. Las tres prioridades se pueden ordenar según lo que pese más en esa compra.

Dieta y alérgenos se declaran dentro del mismo recorrido. La aplicación distingue entre una condición que sí se puede identificar y una que no se puede confirmar. Así, la falta de un dato no se interpreta por sí sola como señal de que el producto es apto.

La información llega por etapas. Primero se muestra el resultado, para facilitar la decisión; después, el detalle, si hace falta. Quien solo quiere elegir entre dos productos no tiene que leer toda la parte técnica para utilizar la aplicación.

La decisión tampoco termina en la ficha. Lo elegido pasa al carrito, donde se pueden revisar las alertas y observar cómo se reparte el conjunto.

Este recorrido está implementado y puede seguirse sin conocer la fórmula que genera el resultado. Sin embargo, no se realizó una prueba formal con compradores ni una medición de satisfacción. Por ello, el proyecto puede describir y demostrar el flujo de uso, pero no afirmar que su usabilidad haya sido validada con personas reales.
